\chapter{PENUTUP}
\label{chap:penutup}

\section{Kesimpulan}
\label{sec:kesimpulan}

Berdasarkan hasil implementasi dan pengujian yang telah dilakukan, maka dapat
diambil beberapa kesimpulan sebagai berikut:

\begin{enumerate}[nolistsep]
  \item Pembuatan platform marketplace tiket konser berbasis smart contract berhasil
        dilakukan menggunakan kontrak ERC-721 yang dikembangkan dalam bahasa Solidity
        dan di-deploy pada jaringan Layer-2 Ethereum, yaitu Polygon zkEVM. Kontrak ini
        mampu mengelola proses minting, validasi penggunaan tiket, dan refund secara
        terdesentralisasi.

  \item Hasil pengujian menunjukkan bahwa fungsi \texttt{purchaseTicket},
        \texttt{useTicket}, dan \\ \texttt{refundTicket} dapat berjalan sesuai dengan
        logika yang telah ditentukan. Transaksi minting tiket berhasil tercatat pada
        block explorer dengan metadata dan informasi kepemilikan yang sesuai dengan
        hasil simulasi di sisi frontend.

  \item Integrasi antara smart contract, antarmuka pengguna berbasis Next.js, dan
        Firebase untuk otentikasi serta penyimpanan data pembelian memberikan
        pengalaman pengguna yang cukup baik dan menunjukkan potensi penggunaan
        teknologi blockchain dalam sistem distribusi tiket yang lebih transparan, aman,
        dan bebas dari pemalsuan.
\end{enumerate}

\section{Saran}
\label{chap:saran}

Untuk pengembangan lebih lanjut terhadap sistem marketplace tiket konser
berbasis smart contract ini, beberapa saran yang dapat dipertimbangkan antara
lain:

\begin{enumerate}[nolistsep]
  \item Memperbaiki sistem validasi penggunaan tiket agar dapat terintegrasi langsung
        dengan sistem fisik (misalnya menggunakan QR code scanner) untuk kebutuhan gate
        checking saat konser berlangsung.

  \item Menambahkan fitur sekunder seperti pasar jual beli tiket bekas (secondary
        marketplace) yang memungkinkan pengguna untuk menjual kembali tiket yang telah
        mereka beli dengan tetap menjaga kontrol kepemilikan melalui NFT.

  \item Mengembangkan UI/UX yang lebih ramah pengguna dan memperluas dukungan terhadap
        jaringan blockchain lainnya seperti Arbitrum, Base, atau StarkNet untuk
        menyediakan pilihan ekosistem Layer-2 yang lebih fleksibel sesuai kebutuhan
        pengguna.
\end{enumerate}
